\section{An actual adapted etale surface chart on the unchanged normalized model}
\begin{theorem}\label{thm:normalized-actual-adapted-etale-chart}
In characteristic zero, retain the original proper parameter model, original DVR and unchanged fraction field, original generic trait, finite relative residue-basis coefficient slice, and original center residue transcendence degree one. The actual normalized model $N$ has an affine open immersion $j:A\to N$ whose image contains the unchanged original center and an actual closed point of its original residual closure. There is an actual etale map $g:A\to\mathbb A_E^2$ preserving the entire original coefficient structure, for which the original residual ideal pulled back along $j$ is exactly the pullback along $g$ of the first coordinate hyperplane. Both the center and the residual curve are the original ones.
\uses{thm:pointwise-regular-divisor-etale-coordinates,thm:normalized-residual-common-smooth-closed-points,thm:normalized-coefficient-slice-center-dvr,thm:normalized-center-original-residue-field}
\lean{EnvelopeProperty.SchemeNormalizedAdaptedEtaleChart.planeBaseMap, EnvelopeProperty.SchemeNormalizedAdaptedEtaleChart.exists_adapted_etale_chart}
\leanok
\end{theorem}
\begin{proof}
Let $C$ be the actual original residual point closure. Its native function field is the unchanged original residue field of the DVR, with its whole $E$-coefficient action; that field has transcendence degree one over $E$. The native function field of $N$ is the unchanged original fraction field and has transcendence degree two over $E$.

The smooth locus of an integral scheme retains the same native function field by the actual generic stalk pullback of its open immersion. This is a whole coefficient-algebra isomorphism: pullback commutes with every original scalar section and its generic germ. Consequently the smooth locus of $N$ is smooth of relative dimension two, and the smooth locus of $C$ is smooth of relative dimension one. Choose an actual closed point $z$ of $C$ where both $C$ and $N$ are smooth. Its image is closed in $N$. The actual local ring of $C$ at $z$ is regular and has dimension one, by the closed-point dimension formula on its smooth locus and the canonical open-immersion stalk isomorphism.

Work on the actual smooth open of $N$ at this image. Pulling back the original residual ideal carries its stalk to the image of its actual original stalk under the open-immersion stalk isomorphism. Quotienting this same isomorphism identifies the resulting quotient with the actual local ring of $C$ at $z$. The quotient is therefore regular of dimension one. The pointwise regular-divisor coordinate theorem now supplies an actual affine neighborhood of the image of $z$, etale over $\mathbb A_E^2$, with its first coordinate cutting out exactly that original residual ideal. Compose its open immersion with the original smooth-open inclusion. Coordinate construction preserves the whole original $E$-structure map.

The generic point of $C$ maps to the unchanged original center of $N$ and specializes to the image of $z$. The image of the new open immersion contains the latter point, so it also contains that original center. Thus all stated objects belong to the original normalized model. This construction uses the reduced original residual closure, and imposes no multiplicity-one or reducedness condition on the whole original parameter fibre.
\end{proof}
